\originalchapter{保角映射}
\sourceinfo{18.04 Topic 10 全部正文; 18.112 L6--L7为主题背景. 交比、单射及边界方向验证为编者补充}
\originalsection{局部与全局}
全纯 $f$ 在 $a$ 满足 $f'(a)\ne0$ 时, 沿任意两条非零切向量 $v_1,v_2$, 像切向量为 $f'(a)v_1,f'(a)v_2$, 它们商不变, 所以有向夹角保持. 这称为局部保角. 若 $f$ 还一一映某域 onto 另一域, 称为两域的双全纯或保角同构. 逆函数全纯已由局部映射定理证明. 局部保角不保证全局单射, 如指数在整个平面.

非恒定全纯单射必有处处非零导数. 若 $f'(a)=0$, 局部映射的重数 $m\ge2$; 选接近但不等于 $f(a)$ 且避开临界值的 $w$, 得至少两不同原像. 更直接, 局部表达 $f-f(a)=h^m$, $h$ 双全纯, $m$ 次幂不单射.

\originalsection{Möbius 变换}
\begin{definition}
映射 $T(z)=(az+b)/(cz+d)$, $ad-bc\ne0$, 称为 Möbius 变换. 在扩充平面上补 $T(-d/c)=\infty$, $T(\infty)=a/c$ ($c\ne0$); $c=0$ 时无穷远映为无穷远.
\end{definition}
代数求逆给 $T^{-1}(w)=(dw-b)/(-cw+a)$, 所以扩充平面一一对应. 有限正常点 $T'=(ad-bc)/(cz+d)^2\ne0$; 在无穷远及极点使用倒数坐标同样为保角. 矩阵 $\begin{pmatrix}a&b\\c&d\end{pmatrix}$ 的非零数倍给同一映射, 复合对应矩阵相乘.

每个 Möbius 变换由平移、非零复数乘法与倒数 $1/z$ 合成. 若 $c\ne0$, 将式子写成 $a/c+(bc-ad)/(c(cz+d))$ 即得分解. 平移和乘法显然将圆、直线映为圆、直线; 倒数的情况可用统一方程证明.
\begin{proposition}
广义圆是扩充平面中的普通圆或带无穷远点的直线. Möbius 变换将广义圆映为广义圆.
\end{proposition}
\begin{proof}
广义圆可写 $A|z|^2+Bz+\bar B\bar z+C=0$, $A,C\in\R$, $|B|^2-AC>0$. $A=0$ 时是直线, $A\ne0$ 时配方得正半径圆. 代入 $z=1/w$ 并乘 $|w|^2$, 形式变为 $C|w|^2+\bar Bw+B\bar w+A=0$, 仍为同类方程. 其他两种基本变换同理保留此形式, 因而任意复合也保留.
\end{proof}

\begin{definition}
不同的四点的交比取约定 $[z,z_1;z_2,z_3]=((z-z_2)(z_1-z_3))/((z-z_3)(z_1-z_2))$, 无穷远项按极限解释.
\end{definition}
恒等式 $T(z)-T(w)=(ad-bc)(z-w)/((cz+d)(cw+d))$ 表明交比不变. 给定三不同源点及三不同像点, 存在唯一 Möbius 变换实现对应: 先用上述交比把源三点 $z_1,z_2,z_3$ 映为 $1,0,\infty$, 再用目标相应变换的逆. 唯一性可由固定 $0,1,\infty$ 的变换只能为恒等证明.

四不同点在同一广义圆上当且仅当交比为实数. 把其中三点所在广义圆映为实轴, 第四点的交比就是其新坐标, 于是结论成立. 这是以三点确定广义圆的解析表达.
\begin{example}
证明 $T(z)=(z-\ii)/(z+\ii)$ 将上半平面一一映为单位圆盘.
\end{example}
\begin{solution}
对 $z=x+\ii y$, $|z+\ii|^2-|z-\ii|^2=4y$, 所以 $y>0$ 恰对应 $|T(z)|<1$. 逆变换 $z=\ii(1+w)/(1-w)$, 对 $|w|<1$ 有 $\Im z=(1-|w|^2)/|1-w|^2>0$. 因而像和逆像都明确, 并非只看实轴映为圆周便断言内部方向.
\end{solution}

\originalsection{圆的反射与同心化}
实轴反射为 $z\mapsto\bar z$, 单位圆反射为 $z\mapsto1/\bar z$. 一般圆 $|z-a|=r$ 的反射为 $a+r^2/(\bar z-\bar a)$, 中心与无穷远互相对应. 其推导是先平移伸缩到单位圆, 反射, 再变换回来. 点与反射点在同一径向线上且到中心的距离乘积为 $r^2$.

对称点的内在描述是: 所有经过两点的广义圆都与给定广义圆正交. 对实轴, 经过共轭两点的圆心位于实轴, 因而正交. Möbius 变换保持广义圆和角, 所以把这种描述及对称点关系一起保持. 这也是反射的唯一性依据.
\begin{proposition}
两条互不相交的广义圆可由 Möbius 变换映为同心圆.
\end{proposition}
\begin{proof}
先取第一条普通圆上的一点映为无穷远, 使它成为直线. 第二条与该点不相交, 所以映为普通圆, 且与该直线不相交. 再用相似变换把圆变为单位圆, 直线变为 $\Re z=c$, 其中可取 $c>1$. 实数 $u=c+\sqrt{c^2-1}$ 和 $v=c-\sqrt{c^2-1}$ 满足 $uv=1$, $u+v=2c$, 因而同时在单位圆和该直线互为反射. 拉回原域得到两点同时关于原两圆对称.

用 Möbius 变换把这两点映为 $0,\infty$. 两个像广义圆都关于 $0,\infty$ 有对称关系, 因而都必须是中心0的普通圆; 直线的反射不会将有限点0映为无穷远. 两圆于是同心. 如果原本其中一条已经是直线, 直接从相似归一化开始即可.
\end{proof}
这一结论也覆盖一圆包住另一圆且不相交的情形. ``不相交'' 排除了相切; 相切时共同对称点合并, 同心化构造退化.

\originalsection{典型区域}
设 $h>0$, $0<\alpha\le2\pi$. 条带 $0<\Im z<h$ 经 $\ee^{\pi z/h}$ 一一映为上半平面, 再接 $T$ 可映为圆盘. 扇形 $0<\arg z<\alpha$ 经 $z^{\pi/\alpha}$ 映为上半平面, 必须使用辐角在 $(0,\alpha)$ 的分支. 一个多值幂公式, 只有指定域与分支后才是一条映射.
\begin{example}
把上半单位圆盘 $|z|<1,\Im z>0$ 一一映为上半平面.
\end{example}
\begin{solution}
$q=(1+z)/(1-z)$ 满足 $\Re q=(1-|z|^2)/|1-z|^2>0$, $\Im q=2\Im z/|1-z|^2>0$, 所以映为第一象限. 逆式 $(q-1)/(q+1)$ 也把第一象限映回该半盘. 再平方得到 $q^2$, 一一映为上半平面. 两端点 $\pm1$ 映为0和无穷远, 半圆弧与直径分别映为负、正实半轴.
\end{solution}
域 $\Re z<0,0<\Im z<\pi$ 经 $\ee^z$ 映为上半单位圆盘, 再用前例可映为上半平面. 这种组合方式把复杂区域的映射分成每一步可核验的一一变换.

设 $a>0$. Joukowski 变换 $J(z)=z+a^2/z$ 将 $|z|>a$ 映为 $\C\setminus[-2a,2a]$. 单射性来自 $J(z_1)=J(z_2)$ 蕴含 $(z_1-z_2)(1-a^2/(z_1z_2))=0$, 外域不允许 $z_1z_2=a^2$. 对任意不在割线的 $w$, 二次方程 $z^2-wz+a^2=0$ 的两根乘积 $a^2$, 恰有一根在圆外: 如果有根在圆周, 写 $z=a\ee^{\ii\theta}$ 则 $w=2a\cos\theta$ 在割线上; 剩余情形由乘积与连续性判定内外, 或利用倒数配对和上、下半平面分支确定. 在上半圆外域, $\Im J=y(1-a^2/|z|^2)>0$, 所以其像为上半平面. 圆周映为线段且上下两岸分别对应不同边界弧.

\originalsection{边值与流动的转移}
若 $T:\Omega\to\Omega'$ 双全纯, $u$ 为 $\Omega'$ 的调和函数, 则 $u\circ T$ 调和. 对可连续对应的边界弧, 边值一起拉回. 在单位圆盘内, 边值右半圆1、左半圆0的解可用 $w=\ii(1+z)/(1-z)$ 转到上半平面. $z=\pm\ii$ 映为 $\mp1$, 右半圆映为 $(-\infty,-1)\cup(1,\infty)$, 左半圆映为 $(-1,1)$. 因而
\[
u(z)=1-\frac{\arg(w-1)-\arg(w+1)}\pi,
\qquad w=\ii\frac{1+z}{1-z},\quad 0<\arg(w\pm1)<\pi.
\]
直接看边界各区间上的角差可核实数值; 中心 $z=0,w=\ii$ 给 $u=1/2$, 与对称性吻合.

复势同样可拉回: 若 $\Phi(w)$ 表示 $w$ 域的流, 则 $\Psi(z)=\Phi(T(z))$ 在 $z$ 域有流线 $\Im\Phi(T(z))=\text{常数}$. 复速度为 $\Psi'=\Phi'(T)T'$, 所以变换会改变速度尺度, 并非只把原速度向量平移过去. 例如以 $T=z^2$ 把第一象限流动变到上半平面问题, 轴边界和半平面实轴由同一变换关联.

\originalsection{习题与解答}
\begin{exercise}
求将 $0,1,\infty$ 分别映为 $1,\infty,0$ 的 Möbius 变换.
\end{exercise}
\begin{solution}
无穷远映为0要求 $a=0$, 1映为无穷要求 $d=-c$, 0映为1要求 $b=d=-c$. 所以 $T(z)=-1/(z-1)=1/(1-z)$, 行列式非零.
\end{solution}
\begin{exercise}
求点 $1+\ii$ 在圆 $|z-2|=3$ 的反射.
\end{exercise}
\begin{solution}
代入 $2+9/(\bar z-2)$, 得 $2+9/(-1-\ii)=-5/2+9\ii/2$. 两点到中心的距离分别为 $\sqrt2$ 与 $9/\sqrt2$, 乘积9.
\end{solution}
\begin{exercise}
求从 $0<\arg z<\pi/3$ 到单位圆盘的双全纯映射, 并核对每步单射.
\end{exercise}
\begin{solution}
$z^3$ 在此扇形辐角变为 $(0,\pi)$ 且模可取任意正值, 一一映为上半平面. 接 $T$ 得 $(z^3-\ii)/(z^3+\ii)$, 上半平面与圆盘的双射已经验证.
\end{solution}
